\documentclass[journal,twoside,web]{IEEEtran}

\usepackage{cite}
\usepackage{amsmath,amssymb,amsfonts}
\usepackage{algorithmic}
\usepackage{algorithm}
\usepackage{graphicx}
\usepackage{textcomp}
\usepackage{xcolor}
\usepackage{booktabs}
\usepackage{multirow}
\usepackage{subcaption}
\usepackage{hyperref}

\newcommand{\methodname}{UCOD-ADF}
\newcommand{\eg}{\emph{e.g.}}
\newcommand{\ie}{\emph{i.e.}}
\newcommand{\etal}{\emph{et al.}}

% Empty placeholder marker for post-training numbers
\newcommand{\PENDING}{---}

\begin{document}

\title{Autonomous Spatial Decision Making for Unsupervised Camouflaged Object Detection}

\author{Computer Vision Research Laboratory
\thanks{This work explores an autonomous multi-agent decision framework for self-supervised camouflaged object detection.}}

\markboth{IEEE Transactions on Pattern Analysis and Machine Intelligence,~Vol.~XX, No.~X,~2026}%
{Research Laboratory: Autonomous Spatial Decision Making for Unsupervised Camouflaged Object Detection}

\maketitle

%==============================================================================
\begin{abstract}
%==============================================================================
Unsupervised Camouflaged Object Detection (UCOD) aims to segment targets visually blended into surrounding environments without human-annotated ground truth masks. Conventional pseudo-labeling pipelines rely on fixed heuristics and global scalar weighting functions that assume spatially uniform label noise. However, pseudo-label noise in camouflaged scenes is inherently non-stationary, concentrating heavily at complex boundaries, low-contrast visual textures, and small instance regions. In this paper, we present \textbf{\methodname{} (Agentic Decision Framework)}, a novel architecture replacing global scalar supervision with an autonomous multi-agent spatial decision system. Our framework integrates four specialized training agents: (1) a \textbf{Quality Assessment Agent (QAA)} utilizing multi-scale ASPP convolutions and feature gradient magnitude sensing; (2) a \textbf{Region Analysis Agent (RAA)} evaluating foreground-to-background context ring contrast and region topology; (3) a \textbf{Memory Agent (MA)} tracking dual-temporal mean and variance dynamics per training sample; and (4) an \textbf{Attention-Gated Decision Agent (DA)} with Spatial and Channel Attention (CBAM) gating to synthesize spatial supervision maps $W_i(x,y) \in [0,1]^{H \times W}$. The framework is optimized end-to-end using three self-supervised loss terms ($\mathcal{L}_{conf}$, $\mathcal{L}_{region}$, $\mathcal{L}_{temporal}$) while maintaining \textbf{zero inference overhead}. We report full quantitative results on four standard COD benchmarks after 25 training epochs.
\end{abstract}

\begin{IEEEkeywords}
Unsupervised Camouflaged Object Detection, Agentic Decision Framework, Spatial Pseudo-Label Calibration, Multi-Agent Learning, Self-Supervised Learning.
\end{IEEEkeywords}

%==============================================================================
\section{Introduction}
\label{sec:introduction}
%==============================================================================

\IEEEPARstart{C}{amouflaged} Object Detection (COD) involves segmenting targets that visually merge into their surroundings through matching textures, boundary patterns, and illumination~\cite{sinet,zoomnet}. Fully-supervised approaches rely on dense pixel-level ground truth masks that are labor-intensive to produce~\cite{camoteacher}. This has motivated research into Unsupervised COD (UCOD), where models learn from unlabeled images using foundation feature backbones and pseudo-labeling heuristics.

Current self-supervised segmentation pipelines suffer from a structural limitation: they apply global, scalar mixing factors across the entire spatial grid, treating clean object interiors and noisy boundary pixels identically. 

%------------------------------------------------------------------------
\begin{figure*}[t]
\centering
\includegraphics[width=0.95\textwidth]{figures/architecture.png}
\caption{\textbf{Teaser Diagram of the \methodname{} System Concept (Figure 1).} Instead of global scalar weighting, \methodname{} processes DINOv2 feature maps and pseudo-labels through a four-agent Autonomous Decision Framework (QAA, RAA, MA, DA) to produce per-pixel spatial supervision maps $W_i(x,y)$ during training.}
\label{fig:teaser}
\end{figure*}
%------------------------------------------------------------------------

To overcome this limitation, we present \textbf{\methodname{} (Agentic Decision Framework)}, formulating pseudo-label supervision as an autonomous multi-agent spatial decision problem during training. By incorporating four specialized training agents (QAA, RAA, MA, DA), our system continuously reasons about spatial reliability, regional difficulty, context contrast, and temporal stability to compute spatial supervision maps $W_i(x,y) \in [0,1]^{H \times W}$.

The main contributions of this paper are:
\begin{enumerate}
    \item We propose \textbf{\methodname{}}, replacing rigid scalar supervision with a spatial multi-agent decision system.
    \item We design four specialized training agents: Quality Assessment (QAA), Region Analysis (RAA), Dual-Temporal Memory (MA), and Attention-Gated Decision (DA).
    \item We introduce three self-supervised loss terms ($\mathcal{L}_{conf}$, $\mathcal{L}_{region}$, $\mathcal{L}_{temporal}$) to train the agentic decision loop end-to-end.
    \item We prove that \methodname{} operates strictly during training, resulting in \textbf{zero parameter or latency overhead during inference}.
\end{enumerate}

%==============================================================================
\section{Related Work}
\label{sec:related}
%==============================================================================

\subsection{Unsupervised Camouflaged Object Detection}
Early unsupervised localization methods like FOUND~\cite{found} and UCOS-DA~\cite{ucosda} leveraged self-supervised Transformer backbones (DINO) to generate background seed masks. UCOD-DPL~\cite{ucoddpl} introduced dynamic pseudo-label mixing and dual-branch adversarial decoders. EASE~\cite{ease} proposed environment-aware prototype libraries. Our work focuses on spatial supervision calibration via multi-agent decision reasoning.

\subsection{Self-Supervised Teacher-Student Frameworks}
The EMA teacher-student framework originates from Mean Teacher~\cite{meanteacher} and BYOL~\cite{byol}. We extend temporal ensembling through a dedicated Memory Agent maintaining sample-level prediction histories.

%==============================================================================
\section{Methodology}
\label{sec:methodology}
%==============================================================================

%------------------------------------------------------------------------
\begin{figure*}[t]
\centering
\begin{subfigure}{0.48\textwidth}
  \centering
  \includegraphics[width=\linewidth]{figures/pipeline_placeholder.png}
  \caption{Macro System Training Pipeline.}
  \label{fig:macro_pipeline}
\end{subfigure}
\hfill
\begin{subfigure}{0.48\textwidth}
  \centering
  \includegraphics[width=\linewidth]{figures/agents_placeholder.png}
  \caption{Internal Multi-Agent Flow.}
  \label{fig:agent_flow}
\end{subfigure}
\caption{\textbf{Detailed Architectural Overview (Figures 2 \& 3).} (a) The complete \methodname{} macro pipeline combining frozen DINOv2 backbone, Teacher/Student decoders, and loss functions. (b) Internal evidence flow across the Quality Assessment, Region Analysis, Memory, and Decision Agents.}
\label{fig:full_architecture}
\end{figure*}
%------------------------------------------------------------------------

\subsection{System Overview}

Given an unlabeled image $I_i$, a frozen DINOv2 backbone extracts features $F_i \in \mathbb{R}^{C \times H \times W}$. A fixed-strategy branch generates background seed pseudo-label $\hat{P}_i^{fs}$. Parallel Student and EMA Teacher Dual-Branch Adversarial (DBA) decoders generate predictions $\hat{Y}_i^{FG}$ and $\hat{P}_i^t$. The agentic decision framework processes these inputs to construct spatial mixing map $W_i(x,y) \in [0,1]^{H \times W}$.

\subsection{Agent 1: Quality Assessment Agent (QAA)}

QAA estimates dense spatial reliability $R_i(x,y)$ using Atrous Spatial Pyramid Pooling (ASPP) with dilation rates $\{1, 2, 4, 6\}$ across a 6-channel input:
\begin{equation}
R_i = \text{QAA}\left([\hat{P}_i^{fs}, \sigma(\hat{P}_i^t), \sigma(\hat{Y}_i^{FG}), D_i, H_i, \|\nabla F_i\|]\right)
\end{equation}
where $D_i = |\sigma(\hat{P}_i^t) - \hat{P}_i^{fs}|$, entropy $H_i = -p\log p - (1-p)\log(1-p)$ with $p = \sigma(\hat{Y}_i^{FG})$, and $\|\nabla F_i\|$ is feature gradient magnitude sensing visual ambiguity. QAA is calibrated via:
\begin{equation}
\mathcal{L}_{conf} = \text{BCE}\left(R_i(x,y),\ \mathbf{1}\left[D_i(x,y) < 0.3\right]\right)
\end{equation}

\subsection{Agent 2: Region Analysis Agent (RAA)}

RAA evaluates connected component topology and background ring contrast. For each component $c_k^{FG} \in \text{ConnectComponent}(\sigma(\hat{Y}_i^{FG}) > 0.5)$, RAA measures:
\begin{equation}
\text{contrast}_k = \|\bar{F}_{c_k} - \bar{F}_{\text{bg\_ring}}\|_2
\end{equation}
where $\text{bg\_ring}$ is a 5-pixel dilated ring surrounding component $c_k$. Regions are classified into Easy, Ambiguous, and Hard categories, penalizing prediction variance on confident regions via:
\begin{equation}
\mathcal{L}_{region} = \frac{1}{|C^{FG}|}\sum_{k} \text{Var}_{c_k}\left[\sigma(\hat{Y}_i^{FG})\right] \cdot \mathbf{1}[\text{label}_k = \text{Easy}]
\end{equation}

\subsection{Agent 3: Memory Agent (MA)}

MA maintains CPU-side sample-level prediction mean $M_i^{(t)}$ and running variance $V_i^{(t)}$ across epochs:
\begin{align}
M_i^{(t)} &= \alpha_m M_i^{(t-1)} + (1-\alpha_m) \sigma(\hat{Y}_i^{FG,(t)}) \\
V_i^{(t)} &= \alpha_m V_i^{(t-1)} + (1-\alpha_m) \left(\sigma(\hat{Y}_i^{FG,(t)}) - M_i^{(t-1)}\right)^2
\end{align}
yielding stability $S_i^{temp} = 1 - |\sigma(\hat{Y}_i^{FG}) - M_i^{(t-1)}|$ and temporal variance $\sigma^2_{temp} = V_i^{(t)}$, with loss:
\begin{equation}
\mathcal{L}_{temporal} = \left\|\sigma(\hat{Y}_i^{FG,(t)}) - M_i^{(t-1)}\right\|_1
\end{equation}

\subsection{Agent 4: Decision Agent (DA)}

DA fuses evidence $[R_i, \text{region\_map}, S_i^{temp}, \sigma^2_{temp}, t/T]$ via Spatial and Channel Attention Gating (CBAM) to compute spatial mixing map $W_i(x,y)$:
\begin{align}
W_i(x,y) &= \sigma\left(\text{MixHead}\left(\text{CBAM}(\mathbf{x})\right)\right) \\
LT_i(x,y) &= \sigma\left(\text{LTHead}\left(\text{CBAM}(\mathbf{x})\right)\right)
\end{align}

The final spatial pseudo-label is blended per pixel:
\begin{equation}
P_i(x,y) = W_i(x,y) \cdot \sigma(\hat{P}_i^t) + (1 - W_i(x,y)) \cdot \hat{P}_i^{fs}
\end{equation}

\subsection{Total Loss Objective}

\begin{equation}
\mathcal{L}_{total} = \mathcal{L}_{seg} + \mathcal{L}_\bot + 0.5 \mathcal{L}_{conf} + 0.1 \mathcal{L}_{region} + 0.3 \mathcal{L}_{temporal}
\end{equation}

%==============================================================================
\section{Training Pipeline}
\label{sec:pipeline}
%==============================================================================

\begin{algorithm}[h]
\caption{\methodname{} Training Procedure}
\label{alg:training}
\begin{algorithmic}[1]
\REQUIRE Dataset $\{I_i\}_{i=1}^N$, max epochs $T$
\STATE Initialize student DBA, teacher DBA (EMA), QAA, RAA, MA, DA
\FOR{$t = 1$ to $T$}
    \FOR{each batch $(I_i, \hat{P}_i^{fs})$}
        \STATE $F_i \leftarrow \text{DINOv2}(I_i)$ \hfill \textit{// Frozen backbone}
        \STATE $\hat{P}_i^t \leftarrow \text{Teacher}(F_i)$ \hfill \textit{// No gradient}
        \STATE $\hat{Y}_i^{FG}, \hat{Y}_i^{BG}, \mathcal{L}_\bot \leftarrow \text{Student}(F_i)$
        \STATE $R_i \leftarrow \text{QAA}(\hat{P}_i^{fs}, \hat{P}_i^t, \hat{Y}_i^{FG}, F_i)$ \hfill \textit{// Agent 1}
        \STATE region\_map $\leftarrow$ RAA$(\hat{Y}_i^{FG}, F_i)$ \hfill \textit{// Agent 2}
        \STATE $S_i^{temp}, \sigma^2_{temp} \leftarrow$ MA.get\_stability$(i, \sigma(\hat{Y}_i^{FG}))$ \hfill \textit{// Agent 3}
        \STATE $W_i, LT_i \leftarrow$ DA$(R_i, \text{region}, S_i^{temp}, \sigma^2_{temp}, t/T)$ \hfill \textit{// Agent 4}
        \STATE $P_i \leftarrow W_i \cdot \sigma(\hat{P}_i^t) + (1-W_i) \cdot \hat{P}_i^{fs}$ \hfill \textit{// Spatial Blend}
        \STATE Compute $\mathcal{L}_{total}$, backpropagate
        \STATE Update Teacher via EMA: $\theta_{teacher} \leftarrow \eta \theta_{teacher} + (1-\eta) \theta_{student}$
        \STATE MA.update$(i, \sigma(\hat{Y}_i^{FG}))$ \hfill \textit{// Memory update}
    \ENDFOR
\ENDFOR
\end{algorithmic}
\end{algorithm}

%==============================================================================
\section{Experiments \& Benchmark Setup}
\label{sec:experiments}
%==============================================================================

\subsection{Experimental Settings}
\textbf{Datasets:} CAMO-Train (1000) + COD10K-Train (3040) without ground-truth labels. Evaluation: CHAMELEON (76), CAMO-Test (250), COD10K-Test (2026), NC4K (4121).

\textbf{Metrics:} S-measure ($S_m$), weighted F-measure ($F_\beta^w$), mean F-measure ($F_\beta^m$), E-measure ($E_\phi$), Mean Absolute Error ($M$), accuracy (ACC), and mean Intersection-over-Union (mIoU).

\subsection{Results at Epoch 25}

\begin{table*}[t]
\centering
\caption{\methodname{} quantitative results across all four test benchmarks, checkpoint at epoch 25.}
\label{tab:epoch25_results}
\resizebox{\textwidth}{!}{%
\begin{tabular}{l|ccccccccc}
\toprule
Dataset & ACC$\uparrow$ & mIoU$\uparrow$ & $E_\phi^{max}\uparrow$ & $E_\phi^{mean}\uparrow$ & $F_\beta^{max}\uparrow$ & $F_\beta^{mean}\uparrow$ & $S_m\uparrow$ & $M\downarrow$ & $F_\beta^w\uparrow$ \\
\midrule
CHAMELEON & .921 & .5788 & .8169 & .8147 & .6408 & .6389 & .734 & .079 & .6162 \\
TE-CAMO & .8703 & .5525 & .7657 & .7636 & .6231 & .6215 & .6962 & .1297 & .5962 \\
TE-COD10K & .92 & .481 & .7465 & .7446 & .5442 & .5425 & .683 & .08 & .5202 \\
NC4K & .9138 & .5809 & .8071 & .8049 & .6455 & .6437 & .7295 & .0862 & .6255 \\
\bottomrule
\end{tabular}%
}
\end{table*}

\subsection{Comparison of Framework Components}

To isolate the effects of the strategy branch, teacher-student branches, and the agentic decision framework, we evaluate three variants under identical conditions (same DINOv2 backbone, same data splits, and same evaluation pipeline):
(1) \textbf{Pure}: The raw baseline without the untrainable Strategy branch or Teacher-Student DBA decoders (evaluated at epoch 25).
(2) \textbf{UCOD-DPL (no agentic)}: The reproduced baseline combining the untrainable Strategy branch (producing $\hat{P}_i^{fs}$) with EMA Teacher/Student DBA decoders, but without the QAA/RAA/MA/DA agents.
(3) \textbf{\methodname{} (with agentic)}: Our proposed model incorporating the multi-agent decision framework on top of the Strategy and Teacher-Student mechanisms.

Table~\ref{tab:agentic_vs_baseline} reports the detailed quantitative results for all three variants.

\begin{table*}[t]
\centering
\caption{Quantitative comparison of the Pure baseline (no strategy, no teacher-student), UCOD-DPL baseline (strategy + teacher-student, no agentic), and \methodname{} (with agentic decision framework) under identical settings.}
\label{tab:agentic_vs_baseline}
\resizebox{\textwidth}{!}{%
\begin{tabular}{ll|ccccccccc}
\toprule
Dataset & Variant & ACC$\uparrow$ & mIoU$\uparrow$ & $E_\phi^{max}\uparrow$ & $E_\phi^{mean}\uparrow$ & $F_\beta^{max}\uparrow$ & $F_\beta^{mean}\uparrow$ & $S_m\uparrow$ & $M\downarrow$ & $F_\beta^w\uparrow$ \\
\midrule
\multirow{3}{*}{CHAMELEON} & Pure (no strategy, no agentic) & .8833 & .5177 & .7411 & .7392 & .5753 & .5737 & .6851 & .1167 & .5476 \\
 & UCOD-DPL (no agentic) & .9687 & .7537 & .9343 & .9316 & .8415 & .8389 & .8647 & .0313 & .8258 \\
 & \methodname{} (with agentic) & .921 & .5788 & .8169 & .8147 & .6408 & .6389 & .734 & .079 & .6162 \\
\midrule
\multirow{3}{*}{TE-CAMO} & Pure (no strategy, no agentic) & .7661 & .4552 & .6282 & .6267 & .5131 & .5120 & .5926 & .2339 & .4821 \\
 & UCOD-DPL (no agentic) & .9231 & .649 & .8638 & .8614 & .7827 & .7805 & .7918 & .0769 & .7471 \\
 & \methodname{} (with agentic) & .8703 & .5525 & .7657 & .7636 & .6231 & .6215 & .6962 & .1297 & .5962 \\
\midrule
\multirow{3}{*}{TE-COD10K} & Pure (no strategy, no agentic) & .8472 & .3784 & .624 & .6225 & .4305 & .4293 & .596 & .1528 & .4024 \\
 & UCOD-DPL (no agentic) & .9699 & .6796 & .9188 & .9162 & .7831 & .7805 & .8342 & .0301 & .764 \\
 & \methodname{} (with agentic) & .92 & .481 & .7465 & .7446 & .5442 & .5425 & .683 & .08 & .5202 \\
\midrule
\multirow{3}{*}{NC4K} & Pure (no strategy, no agentic) & .8448 & .4742 & .6858 & .6841 & .53 & .5287 & .642 & .1552 & .5026 \\
 & UCOD-DPL (no agentic) & .9573 & .7321 & .9248 & .9222 & .8383 & .8357 & .8489 & .0427 & .8171 \\
 & \methodname{} (with agentic) & .9138 & .5809 & .8071 & .8049 & .6455 & .6437 & .7295 & .0862 & .6255 \\
\bottomrule
\end{tabular}%
}
\end{table*}

Across all benchmarks, we observe a clear performance ladder. The Strategy and Teacher-Student (EMA pseudo-labeling) mechanism is the primary driver of performance, raising $S_m$ from $\sim$0.60 in the Pure baseline to $\sim$0.85 in UCOD-DPL. Our agentic variant, \methodname{}, sits between the Pure baseline and UCOD-DPL (yielding $S_m$ values of .734, .6962, .683, and .7295 across CHAMELEON, TE-CAMO, TE-COD10K, and NC4K, respectively). While the agentic decision framework improves results relative to the raw Pure baseline, it regresses performance compared to UCOD-DPL alone. This suggests that the current multi-agent training dynamics might interfere with the convergence of the EMA teacher-student optimization loop, highlighting the need for further tuning of loss weights ($\mathcal{L}_{conf}$, $\mathcal{L}_{region}$, and $\mathcal{L}_{temporal}$) or learning schedules.

\subsection{Comparison with State-of-the-Art}

\begin{table*}[t]
\centering
\caption{Comparison with state-of-the-art methods across four COD benchmarks.}
\label{tab:sota}
\resizebox{\textwidth}{!}{%
\begin{tabular}{l|ccccc|ccccc|ccccc|ccccc}
\toprule
\multirow{2}{*}{Method} & \multicolumn{5}{c|}{CHAMELEON (76)} & \multicolumn{5}{c|}{CAMO-Test (250)} & \multicolumn{5}{c|}{COD10K-Test (2026)} & \multicolumn{5}{c}{NC4K (4121)} \\
& $S_m\uparrow$ & $F_\beta^w\uparrow$ & $F_\beta^m\uparrow$ & $E_\phi\uparrow$ & $M\downarrow$
& $S_m\uparrow$ & $F_\beta^w\uparrow$ & $F_\beta^m\uparrow$ & $E_\phi\uparrow$ & $M\downarrow$
& $S_m\uparrow$ & $F_\beta^w\uparrow$ & $F_\beta^m\uparrow$ & $E_\phi\uparrow$ & $M\downarrow$
& $S_m\uparrow$ & $F_\beta^w\uparrow$ & $F_\beta^m\uparrow$ & $E_\phi\uparrow$ & $M\downarrow$ \\
\midrule
\multicolumn{21}{l}{\textit{Fully-Supervised Methods}} \\
HitNet'23 & .921 & .897 & .900 & .972 & .019 & .849 & .809 & .831 & .906 & .055 & .871 & .806 & .823 & .935 & .023 & .875 & .834 & .853 & .926 & .037 \\
BiRefNet'24 & .929 & .911 & .922 & .968 & .016 & .932 & .914 & .922 & .974 & .015 & .913 & .874 & .888 & .960 & .014 & .914 & .894 & .909 & .953 & .023 \\
\midrule
\multicolumn{21}{l}{\textit{Unsupervised Methods}} \\
FOUND'23 & .829 & .757 & .781 & .911 & .040 & .770 & .704 & .740 & .849 & .090 & .767 & .641 & .668 & .847 & .045 & .816 & .756 & .783 & .893 & .052 \\
UCOS-DA'23 & .750 & .639 & .666 & .808 & .091 & .702 & .604 & .633 & .751 & .148 & .655 & .467 & .495 & .687 & .120 & .731 & .617 & .644 & .785 & .103 \\
UCOD-DPL'25 & .864 & .825 & .838 & .931 & .031 & .793 & .747 & .779 & .862 & .077 & .834 & .763 & .779 & .916 & .031 & .850 & .818 & .835 & .923 & .043 \\
\midrule
\textbf{\methodname{} (Ours)} & .734 & .6162 & .6389 & .8147 & .079 & .6962 & .5962 & .6215 & .7636 & .1297 & .683 & .5202 & .5425 & .7446 & .08 & .7295 & .6255 & .6437 & .8049 & .0862 \\
\bottomrule
\end{tabular}%
}
\end{table}

\subsection{Ablation Study Matrix}

\begin{table}[h]
\centering
\caption{Ablation study matrix evaluating individual agent modules on COD10K-Test. Metric values will be populated after training.}
\label{tab:ablation}
\begin{tabular}{cccc|ccc}
\toprule
QAA & RAA & MA & DA & $S_m\uparrow$ & $F_\beta^w\uparrow$ & $M\downarrow$ \\
\midrule
--- & --- & --- & Baseline & .834 & .763 & .031 \\
$\checkmark$ & --- & --- & Scalar & \PENDING & \PENDING & \PENDING \\
$\checkmark$ & $\checkmark$ & --- & Dense & \PENDING & \PENDING & \PENDING \\
$\checkmark$ & $\checkmark$ & $\checkmark$ & CBAM & .683 & .5202 & .08 \\
\bottomrule
\end{tabular}
\end{table}

%==============================================================================
\section{Complexity Analysis}
\label{sec:complexity}
%==============================================================================

\begin{table}[h]
\centering
\caption{Parameter and Memory Overhead of \methodname{}.}
\label{tab:complexity}
\begin{tabular}{lrrll}
\toprule
Component & Params & FLOPs & Memory & GPU \\
\midrule
QAA & $\sim$48K & $O(420HW)$ & Negligible & Yes \\
RAA & 0 & $O(N_r C)$ & Negligible & No \\
Memory Agent & 0 & $O(HW)$ & 150MB (CPU) & No \\
Decision Agent & $\sim$30K & $O(280HW)$ & Negligible & Yes \\
\midrule
\textbf{Total Multi-Agent} & \textbf{$\sim$78K} & $O(700HW)$ & \textbf{150MB} & --- \\
\bottomrule
\end{tabular}
\end{table}

\textbf{Zero Inference Latency:} Because QAA, RAA, MA, and DA operate exclusively during training, the deployed model at inference time contains zero multi-agent overhead.

%==============================================================================
\section{Conclusion}
\label{sec:conclusion}
%==============================================================================

We introduced \methodname{}, an autonomous multi-agent decision framework for Unsupervised Camouflaged Object Detection. By replacing global scalar mixing with spatial evidence synthesis across Quality Assessment (QAA), Region Analysis (RAA), Dual Temporal Memory (MA), and Attention-Gated Decision (DA) agents, \methodname{} achieves spatial supervision calibration with zero inference overhead. At epoch 25, \methodname{} attains an $S_m$ of .734 on CHAMELEON, .6962 on CAMO-Test, .683 on COD10K-Test, and .7295 on NC4K; current results trail the UCOD-DPL baseline, and further investigation into training convergence and hyperparameter tuning is ongoing.

\bibliographystyle{IEEEtran}
\begin{thebibliography}{10}

\bibitem{ucoddpl}
W.~Yan, L.~Chen, H.~Ko, S.~Zhang, Y.~Zhang, and L.~Cao,
``{UCOD-DPL}: Unsupervised Camouflaged Object Detection via Dynamic Pseudo-label Learning,''
in \emph{Proc. IEEE/CVF Conf. Comput. Vis. Pattern Recognit. (CVPR)}, 2025.

\bibitem{ease}
X.~Hao \emph{et~al.},
``Shift the Lens: Environment-Aware Unsupervised Camouflaged Object Detection,''
in \emph{Proc. IEEE/CVF Conf. Comput. Vis. Pattern Recognit. (CVPR)}, 2025.

\bibitem{sinet}
D.~Fan, G.~Ji, G.~Sun, M.~Cheng, J.~Shen, and L.~Shao,
``Camouflaged Object Detection,''
in \emph{Proc. IEEE/CVF Conf. Comput. Vis. Pattern Recognit. (CVPR)}, 2020.

\bibitem{zoomnet}
Y.~Pang, X.~Zhao, T.-Z. Xiang, L.~Zhang, and H.~Lu,
``Zoom in and out: A mixed-scale triplet network for camouflaged object detection,''
in \emph{Proc. IEEE/CVF Conf. Comput. Vis. Pattern Recognit. (CVPR)}, 2022.

\bibitem{camoteacher}
X.~Lai \emph{et~al.},
``CamoTeacher: Dual-Rotation Consistency Learning for Semi-Supervised Camouflaged Object Detection,''
in \emph{Proc. Eur. Conf. Conf. Comput. Vis. (ECCV)}, 2024.

\bibitem{found}
O.~Sim{\'e}oni \emph{et~al.},
``Unsupervised Object Localization: Observing the Background to Discover Objects,''
in \emph{Proc. IEEE/CVF Conf. Comput. Vis. Pattern Recognit. (CVPR)}, 2023.

\bibitem{ucosda}
O.~Sim{\'e}oni \emph{et~al.},
``Unsupervised Object Localization with Self-supervised Transformers,''
2023.

\bibitem{dino}
M.~Caron \emph{et~al.},
``Emerging Properties in Self-Supervised Vision Transformers,''
in \emph{Proc. IEEE/CVF Int. Conf. Comput. Vis. (ICCV)}, 2021.

\bibitem{byol}
J.-B. Grill \emph{et~al.},
``Bootstrap Your Own Latent: A New Approach to Self-Supervised Learning,''
in \emph{Proc. Adv. Neural Inf. Process. Syst. (NeurIPS)}, 2020.

\bibitem{meanteacher}
A.~Tarvainen and H.~Valpola,
``Mean Teachers are Better Role Models,''
in \emph{Proc. Adv. Neural Inf. Process. Syst. (NeurIPS)}, 2017.

\end{thebibliography}

\end{document}
